% P19-E3: complete two-study scope of Gigerenzer and Hoffrage (1995).
\chapter{Natural Frequencies and Bayesian Reasoning}
\label{sec:p19-natural-frequencies-bayesian-reasoning}

Bayes' formula relates probabilities, but it does not prescribe how a problem
must display its information. A nested count of relevant cases can make a
conditional proportion available through one addition and one division.
Separate percentages may require reconstructing those cases first. Whether
people exploit this difference is an empirical question, distinct from the
validity of Bayes' formula.

This chapter treats the complete two-study design of Gigerenzer and Hoffrage's
1995 paper, \emph{How to Improve Bayesian Reasoning Without Instruction:
Frequency Formats} \cite{GigerenzerHoffrage1995FrequencyFormats}. Its scope is
an elementary inference between two exhaustive hypotheses from one binary
datum. The paper's historical medical and social examples are stipulated
reasoning tasks. Their numbers are not present-day prevalence estimates,
clinical advice, or claims about the social groups named in the stimuli.
The evidence record follows \cref{def:p19-evidence-record}.

\section{A common Bayesian benchmark for different displays}
\label{sec:p19-frequency-benchmark}

\begin{hbdefinition}{def:p19-frequency-menus}{Probability, nested frequencies and information menus}
Let \(H,H^c\) be two hypotheses and \(D\) an evidence event in the finite
probability setting of \cref{thm:p15-finite-bayes}. For \(0<\pi<1\), write
\[
 \pi=\mathbb P(H),\quad u=\mathbb P(D\mid H),\quad
 v=\mathbb P(D\mid H^c),\quad
 x=\pi u,\quad z=(1-\pi)v,\quad s=x+z.
\]
The target is \(\theta=\mathbb P(H\mid D)\), defined when \(s>0\).
The source's \emph{standard menu} supplies \(\pi,u,v\); its \emph{short
menu} supplies \(\mathbb P(H\cap D)=x\) and \(\mathbb P(D)=s\); its
\emph{hybrid menu} supplies \(\pi,u,s\). A menu describes which pieces
are given, while a format describes how those pieces are represented.

For a coherent frequency record let \(a,b,c,d\) be nonnegative integers,
\(N=a+b+c+d>0\), with rows \(H,H^c\) and columns \(D,D^c\):
\[
\begin{array}{c|cc|c}
 &D&D^c&\text{row total}\\\hline
H&a&b&a+b\\
H^c&c&d&c+d\\\hline
\text{column total}&a+c&b+d&N
\end{array}
\]
These are nested counts from one common reference collection. Their
\emph{natural-frequency format} retains subgroup totals on one common scale,
which need not be equal,
as in ``8 of the 10 cases in the first group'' and ``95 of the 990 in
the other.'' The source calls sequential acquisition without fixing marginal
group totals \emph{natural sampling}. A displayed final tally need not
itself have been collected sequentially in the experiment.

By \emph{separately normalized frequencies} we mean conditional frequencies
rescaled to a common denominator within each group, for example 800 per
1,000 \(H\)-cases and 96 per 1,000 \(H^c\)-cases. Their equal written
denominators do not say the groups are equally prevalent. The source's
\emph{relative-frequency} condition uses percentages of groups, retaining
the numerical rates of a probability display. It is distinct from its
natural-frequency condition \cite[pp.~686--689]{GigerenzerHoffrage1995FrequencyFormats}.
\end{hbdefinition}

\begin{hbtheorem}{thm:p19-frequency-bayes-bridge}{proved=annex:p19.frequency-bayes-bridge}{The benchmark and its frequency representation}
With the probability definitions above and \(s>0\),
\[
 \theta=\frac{\pi u}{\pi u+(1-\pi)v}=\frac{x}{s}.
\]
The short and hybrid menus give \(x/s\) and \(\pi u/s\), respectively,
provided their supplied numbers describe this same joint law.
For the uniform distribution on a coherent collection of \(N\) cases,
\[
 \theta=\frac{a}{a+c}\qquad(a+c>0).
\]
This remains defined when one hypothesis row is empty; no rate conditional
on that empty row need be formed. If \(0<a+b<N\), setting
\(\pi=(a+b)/N\), \(u=a/(a+b)\), \(v=c/(c+d)\) gives exactly the same
answer. Every rational triple \((\pi,u,v)\) in the stated domain admits
such an exact integer representation after a common scaling.

Multiplying all four counts by the same positive integer preserves
\(\theta\). Normalizing the two rows separately generally does not:
for \(u,v>0\), using \(u/(u+v)\) gives the same answer as Bayes' formula
exactly when \(\pi=1/2\).
\end{hbtheorem}

The count formula does not discard the prior contribution: it is already
encoded in the relative sizes of \(a\) and \(c\). Supplying hit and
false-positive \emph{rates} requires the group weights; supplying coherent
hit and false-positive \emph{counts} already incorporates them. A ratio
from observed records is an empirical conditional proportion. Using it to
predict a new person's outcome additionally requires a justified relation
between those records and that person. The algebra alone supplies none.

If \(a+c=0\), no evidence-positive case exists in the reference record and
the conditional proportion is undefined. A null prior hypothesis stays null
after positive-probability evidence by \cref{thm:p15-finite-bayes}.
An exact finite uniform representation has rational probabilities; irrational
inputs require approximation rather than an assertion of exact integer counts.
The experiments did not measure the quality of such approximations or infer
a posterior distribution for an unknown population parameter.
The source distinguishes an actual sample size from a generic ``in every
1,000'' rate. Only the former states how many cases were observed. Counts
with that additional information can support an uncertainty analysis under
a specified sampling model and, for a Bayesian parameter analysis, a prior;
neither is supplied merely by choosing a convenient display denominator.
The source explicitly leaves its result about parameter uncertainty outside
the experiments' single-point-estimation scope
\cite[pp.~687--689]{GigerenzerHoffrage1995FrequencyFormats}.

\begin{hbexample}{def:p19-mammography-formats}{The four mammography versions and their rounding}
In this historical task, \(H\) is breast cancer for a woman aged forty
participating in routine screening, and \(D\) is a positive mammogram.
The source's standard probability task supplies
\(\pi=0.01,u=0.80,v=0.096\). Its exact benchmark is
\[
 \frac{0.008}{0.008+0.09504}=\frac{25}{322}
 \approx0.07763975.
\]
An exactly matched integer table has \(N=12500\), with
\((a,b,c,d)=(100,25,1188,11187)\), yielding \(100/1288=25/322\).

The source's standard frequency display instead gives 10 of 1,000 in \(H\),
8 of these 10 in \(D\), and 95 of the other 990 in \(D\). It yields
\[
 (a,b,c,d)=(8,2,95,895),\qquad \theta=8/103\approx0.07766990.
\]
Its short probability display gives \(x=0.008,s=0.103\), and its short
frequency display gives 8 joint cases and 103 evidence-positive cases per
1,000. Both give \(8/103\). The slight difference from \(25/322\) is
rounding: \(95/990\ne0.096\). All round to about 7.8\%, but their
printed inputs are not exactly equivalent
\cite[pp.~685--688, table~1]{GigerenzerHoffrage1995FrequencyFormats}.

For contrast, pooling separately normalized rows with 800 positive results
per 1,000 \(H\)-cases and 96 per 1,000 \(H^c\)-cases gives
\(800/(800+96)=25/28\approx89.29\%\). This silently substitutes a
50\% prior for 1\%. A change from percentages to ``per 1,000'' therefore
does not by itself create a valid natural-frequency calculation.
\end{hbexample}

\section{Algorithms, shortcuts and the response criterion}
\label{sec:p19-frequency-algorithms}

For the explicitly displayed arithmetic, compute \(x=\pi u\) and
\(z=(1-\pi)v\), add them and divide. Given \(a,c\), add and divide;
given \(a,a+c\), only divide. Given the hybrid menu, multiply and divide.
These are counts of operations in particular formulas, not a lower bound
for every possible algorithm or a theorem about mental effort. Converting
a probability display into counts also requires work. The source's short
frequency menu already gives the denominator, saving the addition needed
in its standard frequency menu
\cite[pp.~687--689, 694]{GigerenzerHoffrage1995FrequencyFormats}.

The paper also describes a participant's beam construction: represent the
whole population by a unit length, split it into lengths \(\pi,1-\pi\),
retain proportions \(u,v\) of the respective pieces, then compare the
retained \(H\)-piece with their sum. Its lengths are exactly \(x,z\),
so the comparison is \(x/(x+z)\). The drawing is one observed solution
method, not evidence that every participant used pictorial representations
\cite[pp.~689--690, Fig.~2]{GigerenzerHoffrage1995FrequencyFormats}.

\begin{hblemma}{thm:p19-bayesian-shortcut-errors}{proved=annex:p19.bayesian-shortcut-errors}{Exact errors for the reported shortcut families}
Let \(0<\pi<1\), \(u,v\in[0,1]\), \(x=\pi u\),
\(z=(1-\pi)v\), and \(x+z>0\). Whenever the displayed denominators
are positive, define the rare-event and high-hit-rate approximations
\(R=x/(x+v)\) and \(T=\pi/(\pi+z)\). Then
\[
 \theta-R=\frac{x\pi v}{(x+z)(x+v)},\qquad
 T-\theta=\frac{\pi z(1-u)}{(\pi+z)(x+z)}.
\]
For the comparison approximations,
\[
 \frac{x}{z}-\theta=\frac{x^2}{z(x+z)}\quad(z>0),\qquad
 \theta-\left(1-\frac{z}{x}\right)=\frac{z^2}{x(x+z)}\quad(x>0).
\]
If \(x/z\leq\varepsilon\), the first error is at most
\(\varepsilon^2\); if \(z/x\leq\varepsilon\), the second is at most
\(\varepsilon^2\). These statements quantify the relevant small-ratio
condition. They do not make either approximation exact.
\end{hblemma}

The source's combined shortcut uses \(\pi/v\), with \(v>0\), when \(H\) is rare, \(u\)
is near one and false-positive joint cases dominate; the reverse-dominance
version uses \(1-v/\pi\). Such formulas are approximate and can leave
\([0,1]\) outside suitable conditions. They must be assessed against the
benchmark, not classified as Bayesian merely because they are short
\cite[pp.~690--691]{GigerenzerHoffrage1995FrequencyFormats}.

The source's four empirical predictions were more Bayesian-coded solutions
with frequency than probability formats; more with short than standard
probability menus; equal proportions across frequency menus; and similarly
low proportions for probability and relative-frequency formats. Study~2
extended the probability-menu prediction by putting the hybrid between
standard and short. Moving from an operation count to any such prediction
requires a psychological hypothesis about how people use the display; it
is not a deductive consequence of Bayes' theorem
\cite[pp.~691--692, 695]{GigerenzerHoffrage1995FrequencyFormats}.

\paragraph{The actual coding rule.}
Study~1 required both a compatible numerical answer and a written process
showing a Bayesian calculation, pictorial analogue, or qualifying shortcut.
The source allowed rounding an exact percentage either down or up to the
adjacent whole percentage. It accepted a correct unevaluated fraction and
ignored a division error when the preceding equation was Bayesian; the goal
was not to score elementary arithmetic. Consequently this is not simply a
count of exactly correct final numerals. A shortcut whose result fell outside
the allowed outcome range was assigned elsewhere, even if its general form
was recognized \cite[p.~692]{GigerenzerHoffrage1995FrequencyFormats}.

Participants were asked to write their reasoning, using calculations,
diagrams or other marks unless guessing, and to say whether each answer was
calculated or guessed. Unclear calculation records prompted a few interviews
after all tasks; questioning stopped if participants could not immediately
interpret their notes. These protocols supply additional evidence about a
solution, but are not a complete direct observation of cognition. The paper
does not report a numerical inter-coder reliability coefficient.

\section{Study 1: crossed formats and menus}
\label{sec:p19-frequency-study-one}

The sample comprised 60 paid University of Salzburg students from ten
disciplines, predominantly psychology: 21 men and 39 women, median age 21.
None was familiar with Bayes' theorem. Participants worked individually or
in groups of two or three, with two groups of five. Sessions were expected
to last about an hour, but extra time was allowed. Actual mean durations
were 73 minutes (range 25--180) and 53 minutes (30--120).
Recruitment dates, payment amount and a population probability-sampling
procedure are not reported \cite[pp.~692--693]{GigerenzerHoffrage1995FrequencyFormats}.

The two formats, probability and natural frequency, were crossed with
standard and short menus. Participants were randomly assigned to two groups:
one received standard probability and short frequency versions; the other
received standard frequency and short probability versions. Each answered
each of 15 problems twice, once per assigned version, across two sessions
one week apart. Formats and menus were balanced across sessions; versions
of the same problem appeared in different sessions. Two random problem
orders were used within each group. The resulting 30 tasks per person are
repeated observations, not 30 independent participants.

Probability-format questions requested a single-event probability in percent.
Frequency-format questions requested an expected number of \(H\)-cases
among a new representative sample of \(D\)-cases, with both blanks in
an ``out of'' answer. Input format and requested output format therefore
changed together. The study is not a pure manipulation of Arabic numerals
with every other wording and response requirement held fixed.

\paragraph{Complete problem and numerical-input scope.}
The following table preserves the 15 problem topics and the three source
ratios. Each fraction is the printed ``out of'' pair, not an assertion that
all denominators form one exact partition. The final column reproduces the
source's reported benchmark percentage
\cite[p.~693, table~2]{GigerenzerHoffrage1995FrequencyFormats}.
\begin{center}\small
\begin{tabular}{lrrrr}
Problem: hypothesis/evidence&\(H\)&\(D\mid H\)&\(D\mid H^c\)&\(100\theta\)\\\hline
Cancer/positive mammogram&10/1000&8/10&95/990&7.77\\
Prenatal damage/measles&21/10000&10/21&50/10000&16.70\\
Blue cab/witness says blue&15/100&12/15&17/85&41.38\\
AIDS/positive HIV test&100/1000000&100/100&1000/1000000&9.09\\
Heroin addiction/needle mark&10/100000&10/10&190/100000&5.00\\
Pregnancy/positive test&20/1000&19/20&5/980&79.17\\
Car accident/driver drunk&100/10000&55/100&500/9900&9.91\\
Child's posture/heavy books&50/1000&20/50&190/950&9.52\\
School accident/urban home&30/1000&27/30&388/970&6.51\\
Suicide/professor&240/1000000&36/240&120000/1000000&0.03\\
Red ball/star marking&400/500&300/400&25/100&92.31\\
Economics course/career focus&300/1000&210/300&350/700&37.50\\
Active feminist/bank teller&5000/100000&20/5000&2000/95000&0.99\\
Pimp/Rolex watch&50/1000000&40/50&500/1000000&7.41\\
School admission/test result&360/1000&270/360&128/640&67.84
\end{tabular}
\end{center}

The source explicitly rounded some large subgroup denominators for ease of
presentation. Hence its prenatal-damage row has \(H^c\) denominator 10,000
rather than \(10000-21\). Interpreting all three printed pairs as rates
gives 16.695884\%, while treating 10 and 50 as coherent joint counts gives
16.666667\%. Both approximations explain the intended scale, but they
are not identical inputs. The complete matched calculations and other rounded
rows are supplied in \cref{thm:p19-frequency-numerical-reconstructions}.

\paragraph{Main results and actual denominators.}
Table~3 reports 1,774 answers, against \(60\times30=1800\) opportunities;
26 were absent through refusals and missing data. The later discussion
identifies 21 refusals, 17 in the standard probability condition. It does
not justify reallocating the other missing records to a guessed category.
The following numerators and denominators are source counts; the rounded
percentages are those reported in the prose.
\[
\begin{array}{c|rr}
 &\text{standard menu}&\text{short menu}\\\hline
\text{Probability}&69/433\ (16\%)&126/450\ (28\%)\\
\text{Natural frequency}&204/445\ (46\%)&221/446\ (50\%)
\end{array}
\]
The numerator is the number assigned the Bayesian process/outcome code.
Unidentified solutions remain in the denominator. The paper's approximately
50\% result is neither the proportion of people who always reasoned
Bayesianly nor the proportion restricted to identifiable solutions
\cite[pp.~693--695, 699]{GigerenzerHoffrage1995FrequencyFormats}.

To preserve the remaining outcome distribution, the next table gives all
source categories. Columns PS, PQ, FS and FQ denote probability standard,
probability short, frequency standard and frequency short. A zero denotes
a blank count cell in the source table.
\begin{center}\small
\begin{tabular}{lrrrrr}
Source response code&PS&PQ&FS&FQ&Total\\\hline
Bayesian&69&126&204&221&620\\
Joint occurrence&39&97&20&97&253\\
Adjusted joint occurrence&0&64&0&55&119\\
Likelihood picked (source: Fisherian)&67&0&36&0&103\\
Adjusted likelihood picked&32&0&19&0&51\\
Multiply all&0&79&0&12&91\\
Likelihood subtraction&30&0&4&0&34\\
Base rate only&6&0&13&0&19\\
Less frequent algorithms&71&32&60&29&192\\
Not identified&119&52&89&32&292\\\hline
Total&433&450&445&446&1774
\end{tabular}
\end{center}
Joint occurrence selects \(x\); its adjusted variant changes the estimate
by at most five percentage points. Likelihood picking selects \(u\), with
an analogous adjusted variant. Multiply all in the short menu uses \(sx\);
likelihood subtraction uses \(u-v\); base-rate-only uses \(\pi\).
The source's term \emph{Fisherian} refers to treating a likelihood as a
posterior, not to a correct Fisher significance test. Similarly, likelihood
subtraction is not the Neyman--Pearson likelihood-ratio procedure. The source
makes the first qualification explicitly. Each unlisted rare algorithm
occurred in less than 1\% of all solutions; their individual forms are not
specified in the paper \cite[pp.~694--695]{GigerenzerHoffrage1995FrequencyFormats}.

The source reports format differences of 29.9 and 21.6 percentage points
(standard and short menus) and menu differences of 12.1 and 3.7 points
(probability and frequency formats). The residual frequency-menu difference
fits the saved addition in the short menu; it is not exact equality of
performance. Figure~3 plots all 15 items; the text reports that standard
probability yielded the lowest Bayesian-code proportion in each and that
short exceeded standard within probability formats in each. Its specifically
printed short-frequency mammography rate is 43.3\%. The plotted item bars
are not reverse-engineered into unreported exact counts.

For the standard-frequency suicide item, 9 of 30 participants used the
comparison shortcut and 5 used a Bayesian algorithm without a shortcut.
Across 20 observed comparison-shortcut uses, 17 met the strict outcome rule;
the other 3 were within four percentage points. The 20 uses concern repeated
solutions, not a further sample of 20 people
\cite[pp.~693--694]{GigerenzerHoffrage1995FrequencyFormats}.

\section{Study 2: probability and relative frequency across three menus}
\label{sec:p19-frequency-study-two}

Study~2 was performed \emph{before} Study~1, although presented second.
It enrolled 15 paid University of Konstanz students from biology, linguistics,
English, German, philosophy, political science and management: 8 men and
7 women, median age 22, none familiar with Bayes' theorem. They worked in
one group. A prize of 25 deutsche marks was offered for the best performance
\cite[pp.~695--696 and note~4]{GigerenzerHoffrage1995FrequencyFormats}.

There were 24 problems, twelve reused in Study~1 and twelve new; some shared
problems had different numbers across studies. The paper does not print a
complete 24-item stimulus set. Each problem appeared in all three menus,
giving 72 tasks per participant. Half the problems were randomly assigned
to probability wording and half to relative-frequency wording; both used
percentages, and both requested a single-event probability. Thus each person
encountered both formats and all menus. This study contained no natural-
frequency input condition.

Six scheduled one-hour sessions were spread equally over three weeks, with
twelve tasks per session. Menu versions of a problem never occurred in the
same session, and all three menus occurred equally often in each session.
Subject to these constraints, tasks, within-session order, session allocation
and order of the numerical information were randomized as described in the
source. No Bayesian instruction or outcome feedback was given. The title's
``without instruction'' means without instruction in Bayesian solution
methods, not absence of task instructions or incentives.

Unlike Study~1, this study did not require written concurrent protocols.
Participants could write notes spontaneously. After all 72 tasks, each
received copies of six of their completed tasks and answered questions about
the information and reasoning used. This supplies less concurrent process
evidence and adds retrospective-report limitations. The study identified
algorithms for about 67\% of its 1,080 judgments; it should not be described
as having the same complete instructed protocol for every answer as Study~1
\cite[p.~696]{GigerenzerHoffrage1995FrequencyFormats}.

\paragraph{Every reported format and menu condition.}
The Bayesian counts were 60/540 for relative-frequency wording and 66/540
for probability wording. Pooling the two formats as the source does, the
counts were 23/360, 40/360 and 63/360 for standard, hybrid and short menus,
respectively (approximately 6.4\%, 11.1\%, 17.5\%). The paper reports
descriptive similarity across formats and increasing proportions across
menus. It gives no equivalence test, confidence interval or null-hypothesis
test establishing identical format effects.

Table~4's complete counts follow. RF and P denote relative frequency and
probability; S, H and Q denote standard, hybrid and short menus. These are
two marginal breakdowns of the same 1,080 answers, not separate samples
and not a published six-cell joint table.
\begin{center}\small
\begin{tabular}{lrrrrrr}
Source response code&RF&P&S&H&Q&Total\\\hline
Joint occurrence&91&88&46&31&102&179\\
Bayesian&60&66&23&40&63&126\\
Likelihood picked&46&45&41&50&0&91\\
Adjusted likelihood picked&20&29&20&29&0&49\\
Multiply all&11&27&0&3&35&38\\
False-alarm complement&17&20&37&0&0&37\\
Likelihood subtraction&19&9&28&0&0&28\\
Base rate only&14&10&14&10&0&24\\
Total negatives&10&7&0&9&8&17\\
Positives times base rate&7&7&0&14&0&14\\
Positives times hit rate&4&9&0&13&0&13\\
Hit rate minus base rate&6&5&3&8&0&11\\
Less frequent algorithms&60&37&37&34&26&97\\
Not identified&175&181&111&119&126&356\\\hline
Total&540&540&360&360&360&1080
\end{tabular}
\end{center}
The additional formula labels are \(1-v\), \(1-s\), \(s\pi\),
\(su\), and \(u-\pi\), in the corresponding row order. Multiply-all
can use the displayed factors to form \(sx\), including \(s\pi u\)
when the hybrid menu supplies \(\pi,u,s\). Such arithmetic is a coded
response strategy, not a Bayesian posterior identity.

\paragraph{Dependence across versions and over time.}
There were \(24\times15=360\) person-by-problem triples across menus.
The same identified algorithm appeared across all three menus in only 16
triples: 10 Bayesian and 6 joint-occurrence. This does not prove that every
other triple used three different processes; unidentified responses also
limit the comparison. The mean within-triple range of estimates was 40.5
percentage points across all problems, and 41 points for the measles example.
These are ranges of a person's repeated estimates, not standard errors.

Across Study~2's first three versus final three sessions, the reported Bayesian
proportions changed from .04/.07/.14 to .08/.14/.21 for standard/hybrid/short
menus. Study~1 changed from .30 to .38 over its two sessions. These are
practice-associated descriptions without a separately randomized practice
control. The source's temporal consistency measure is the fraction of adjacent
same-format/menu problems given the same algorithm code. Study~1 frequency
consistency rose from .32 to .49 (standard) and .61 to .76 (short); standard
probability changed .22 to .24. Study~2 standard-probability consistency was
.15 then .16; short-probability values averaged across the studies were .40
then .42 \cite[pp.~696--697 and note~6]{GigerenzerHoffrage1995FrequencyFormats}.

\section{Matched reconstructions and limits of the process interpretation}
\label{sec:p19-frequency-reconstructions}

\begin{hbtheorem}{thm:p19-frequency-numerical-reconstructions}{proved=annex:p19.frequency-numerical-reconstructions}{Exact arithmetic for the source's task versions}
The mammography calculations in \cref{def:p19-mammography-formats} follow
from the finite Bayes benchmark. For source pairs \(h/N,a/h,c/m\), assume
positive integers \(N,h,m\), \(0<h<N\), nonnegative integers \(a,c\),
\(a\leq h\), \(c\leq\min(m,N-h)\), and \(a+c>0\).
Their literal rate interpretation yields
\[
 \theta_{\rm rates}=\frac{am}{am+(N-h)c}.
\]
The coherent-leaf interpretation yields \(a/(a+c)\); it agrees when
\(m=N-h\), but need not agree after rounding that denominator.

In the measles task, \(H\) denotes severe prenatal damage and \(D\)
German measles in the mother during early pregnancy. Study~2's standard
menu supplies \(\pi=.005,u=.4,v=.0001\)
and yields \(4000/4199\approx.95260776\). Its hybrid and short menus
supply \(s=.0021\), giving \(.002/.0021=20/21\approx.95238095\).
Thus these printed menus also agree only approximately. For the standard
inputs the rare-event, high-hit-rate and reverse-comparison shortcuts yield
\(20/21\), \(10000/10199\) and \(.95025\), respectively. The combined
reverse shortcut gives \(.98\).
The separate HIV shortcut illustration on p.~691 uses
\(\pi=.002,u=.99,v=.02\), not Study~1's HIV numbers. Its exact
posterior is \(99/1097\approx.09024613\), while \(\pi/v=.10\).
For the standard mammography inputs, the combined shortcut gives
\(5/48\approx.10416667\), compared with \(25/322\).
\end{hbtheorem}

The source's four Appendix cases preserve why outcome and process are
different records. In standard/hybrid/short order their reported percentage
answers were R.: 35/.5/.2; M.: 40/.1/95.7; S.: .2/.04/.04; and
O.: 100/.5/95.3. The cases illustrate selected individuals, not a random
subsample establishing prevalence. R. moved among adjusted likelihood,
base-rate-only and joint-occurrence codes; M. moved from likelihood picking
to multiplying positive rate by hit rate to a Bayesian frequency conversion;
S. used joint occurrence then multiply-all; O. used false-alarm complement,
base-rate-only and a Bayesian ratio
\cite[p.~704]{GigerenzerHoffrage1995FrequencyFormats}.

In M.'s final reasoning, the complement count was 10 out of 210; its correct
percentage is 4.7619\%, giving 95.2381\% for the target, although the
recorded final answer was 95.7\%. The source classified the method Bayesian
under its division-error allowance. S.'s recorded .04\% is likewise retained
as a response: the stated multiply-all factors \(.005\times.0021\times.4\)
actually give .00042\% after conversion to percent, not .04\%. The
annex supplies these computations; an observed arithmetic error is not
silently converted into a correct benchmark.

\paragraph{What is and is not identified by a response code.}
The source checks two broader candidate rules, picking one number and
integrating everything, against menu triples; 12 of 360 were consistent
with the former and 14 with the latter. These counts challenge an account
using either rule unchanged throughout, but do not uniquely identify every
uncoded process. The article also reports across-problem Pearson correlations:
joint-occurrence use with \(s\), .36 and .47 in Studies~1 and~2;
likelihood-picking use with \(\lvert1-\pi/s\rvert\), -.49 and -.59.
Their units are 15 and 24 problems, respectively. No confidence intervals
or tests are supplied for these correlations
\cite[pp.~698--699]{GigerenzerHoffrage1995FrequencyFormats}.

Under the source's algorithm classifications, false-alarm-rate neglect was
reported in 31\% and 33\% of cases, and base-rate neglect in 32\% and
36\%, excluding contributions from less frequent algorithms. The argument
is that different processes can omit the same input and can produce similar
answers. Neither omission nor a unique process follows from a numerical
answer alone \cite[pp.~699--700]{GigerenzerHoffrage1995FrequencyFormats}.

\paragraph{Uncertainty and generalization.}
The main comparisons are reported through counts, percentages and descriptive
differences. The article supplies no confidence intervals or inferential
tests for the principal format/menu contrasts, and the underlying participant
records needed to calculate uncertainty while respecting dependence are not
published with it. A binomial interval treating all 1,774 or 1,080 answers
as independent would add an unsupported model. Missing records, unequal
process reporting across studies and repeated exposure remain material.

Crossing format with menu challenges the proposal that fewer information
pieces alone explain Study~1's result. It does not separately randomize
evolutionary history, familiarity with numbers, internal representations,
response wording, or all possible processing explanations. The authors
explicitly avoid specifying the internal representation and restrict their
demonstrations to binary hypotheses/data and one datum. Their comments on
multiple dependent cues, continuous variables, other biases and teaching
systems are broader arguments or proposals, not additional experiments
performed in this paper \cite[pp.~685, 698--701]{GigerenzerHoffrage1995FrequencyFormats}.

Natural-frequency displays in these tasks often supported the coded Bayesian
method without teaching that method. They did not make every answer Bayesian,
establish a unique evolved mechanism, evaluate patient outcomes, or compare
actual clinical decision policies. A reference user can take the proved
representation bridge and the recorded experimental contrasts separately:
first ensure that the numerical display describes the intended joint law,
then assess whether the relevant users interpret and use it as intended.
